% Generated draft: requires mathematical and source review.
\documentclass[11pt]{article}
\usepackage[margin=25mm]{geometry}
\usepackage{fontspec}
\setmainfont{Latin Modern Roman}
\usepackage{amsmath,amssymb,mathtools,unicode-math}
\setmathfont{Latin Modern Math}
\usepackage{enumitem,needspace,xurl,hyperref}
\setlength{\parindent}{0pt}
\setlength{\parskip}{5pt}
\setlength{\emergencystretch}{3em}
\newcommand{\N}{\mathbb{N}}
\newcommand{\Z}{\mathbb{Z}}
\newcommand{\Q}{\mathbb{Q}}
\newcommand{\R}{\mathbb{R}}
\newcommand{\C}{\mathbb{C}}
\newcommand{\F}{\mathbb{F}}
\newcommand{\A}{\mathbb{A}}
\newcommand{\PP}{\mathbb{P}}
\newcommand{\E}{\mathbb{E}}
\newcommand{\eps}{\varepsilon}
\newcommand{\dd}{\,\mathrm{d}}
\begin{document}
\section*{UnsolvedMath ID: 30001971}
\textbf{Source problem code:} catalogue entry
\subsection{Definitions and mathematical statement}
The supplied record concerns omitted steps in a historical mathematical algorithm. It identifies neither the algorithm nor its author, the omitted steps, or the surviving textual witnesses. Consequently, no specific mathematical objects, hypotheses, domains, boundary conditions, or quantified proposition can be recovered from the supplied evidence.

The source question is:
\begin{quote}
Why did the author of a mathematical algorithm omit some steps: by deliberate choice, or because the surviving text is corrupted?
\end{quote}
Here deliberate omission and corruption of the surviving text are proposed historical explanations. The supplied record gives no criteria for distinguishing them and does not establish that they are mutually exclusive or exhaustive. A precise mathematical statement cannot therefore be supplied without inventing information.

The supplied statement-verification annotation classifies this as a historical, pedagogical, or methodological discussion prompt rather than an unresolved mathematical proposition. Its statement status is \texttt{unrecoverable}, and no separate clean statement is supplied. The registry label \texttt{open} is retained as supplied metadata; it does not establish the existence of a well-defined open mathematical conjecture.
\subsection{Short English statement}
The question asks whether missing steps in an unspecified historical algorithm reflect the author's deliberate choice or corruption of the surviving text. The supplied evidence does not identify the algorithm or the relevant manuscripts, so a precise mathematical proposition cannot be recovered. The original historical question is preserved for review.
\subsection{Sources}
\begin{itemize}
\item Local UnsolvedMath registry, permanent ID \(30001971\), title \emph{Omitted Steps in Historical Mathematical Algorithms}, source problem code \texttt{catalogue entry}, in the Oberwolfach Reports Open Problems collection. The record gives proposed year \(2012\) but identifies no proposer or algorithm author.
\item Supplied source citation: \emph{Explicit Versus Tacit Knowledge in Mathematics} (2012), Oberwolfach Reports, pp.~131--210, DOI \texttt{10.4171/owr/2012/04}. Supplied source URL: \url{https://doi.org/10.4171/owr/2012/04}. The supplied background describes Oberwolfach Reports as documenting mathematical workshops and their open-problem discussions; the statement-verification annotation characterizes this particular item as a historical, pedagogical, or methodological discussion prompt.
\item Supplied literature source, labeled ``original Oberwolfach report'': \url{https://doi.org/10.4171/owr/2012/04}. Literature checked at: \texttt{2026-08-21}. This date and assessment are supplied dated curation, not independently verified. The supplied assessment labels the item \texttt{open} and states that no determinate answer to the historical question about omitted steps and textual corruption is established by the extracted report passage. The background requests independent verification of time-sensitive status and cited claims.
\end{itemize}
\end{document}
